\section{All Soul's Day}

Around the second of November, I would like to avoid naming here the living, or rather to be concerned with them only insofar as they themselves are concerned with men who had already departed from life. A melancholy memory is not a simple dream, and nothing deep down is more useful to those who remain than the strong tenor of those who have left.

But perhaps it would be suitable, while on this subject, to make a distinction. There is the universal cult of the dead, of all the dead, of those who had existed, provided that they had belonged to the human race; and there is, closely related, the particular cult, more reserved, prouder, and, in my sense, more beneficial. That renders to the elite among the dead, those whom the positivists call, a little verbosely, “the great types of humanity”, and the Catholics, more briefly, the “saints”. The first of these cults presents a great drawback; in teaching us to venerate all of defunct humanity, it trains us logically to venerate, en bloc, all of living humanity, that is to say, to make us accept and even venerate the worse faults that it commits even though we recognize them as much in ourselves as in our neighbours. The second cult shows the opposite benefit; by obliging us to hold the dead as our models, it forces us to select from among these scattered people, hence, indirectly, to make a critique of our own characters: by applying our minds to consider those great dead men, it opens us up to the way of personal exaltation and perfection.

The consequence is that human solidarity, to use that term, must belong much less to the crowd of our predecessors, than to the persons of the past who have realized, in a great way, the fine natural traits of man. Those who pass up the opportunity to serve their great memory, pass up an undoubted opportunity to help themselves, to correct themselves, and to improve themselves.

That said, what to think of the cold and sad silence that has greeted these last few weeks—in the newspapers, in the reviews, everywhere finally where they pride themselves on their thought and style—the great news of the discovery of a newly published book of Bossuet? Regarding the second tract on the États d’oraison (“States of Prayer”), if it had taught nothing new of the opinions and sentiments of the illustrious bishop, if it had been only a sketch, a pale appendix to some known works, it still would have been necessary to honour in it, after two hundred years, some unheard contemplative words of a voice which fills the world. This indifference of the critics and the public, had, moreover, no excuse, and the six chapters published the other day in la Quinzaine by the fortunate author of the discovery, Abbot Levesque of Saint‐Sulpice, was worthy of attention and of commentary for the same reason that \textit{Relation sur le quiétisme} (``Relation to quietism''), so fine, so piquant and solid, especially because of the first tract of États d’oraison whose eloquence equals the power of the ideas.

The honour of our times will have been to render to Bossuet his rightful place. We have displaced him from it. First, by literary rancour: reread Paul Albert and the other critics of romanticism, and you will see how much this great man baffles them. It is that, if it is not right, at least it is possible to contend that a Jean Racine or a Jean de La Fontaine have united a marvellous vigour to their supreme perfection: to deceive the inane, or perhaps by innate taste in reserve, these two poets have disguised their strength under elegance and grace. But this force bursts out divinely in Bossuet. I do not know which imagination was more powerful than his or more allied, for example, to that of Homer; he joins to it a sentiment, an almost dramatic passion: nevertheless, order is never weakened in him. It is not necessary to be astonished if true strength is naturally harmonious and there is always something weak or morbid, even vertiginous, in disorder. The most fanciful of men—Plato, Aristophanes—were the most reasonable, the most sensible, and the finest.

After the romantics, Bossuet bored Renan. Pardon the word, let us rather not apply it to the person of Renan, but to the adjuncts, the dependants, if you like, to the low offices of this illustrious writer: there were in him some parts corrupted by Hegelianism. Bossuet's unswerving logic resisted it. I imagine that it would still resist it if, by some miracle, Bossuet returned among us. Everything that the Germans said against the principles of reason is something confused. They always mixed the orders. All the same, in a superior order, our reason would be defective, there remains an order, the human order, where, in spite of all, it exercises itself supremely. That is what the author of La Synthèse subjective discerned with true genius. The influence of Comte has led the wise public back to Bossuet; as for those who don’t have time for reflection, the fortunate intervention of Mr. Brunetière produced around the same time the same happy effect.

So that the most harmonious of men is admired with unanimity.

The silence maintained on the new treasure added to the old treasure by M. Levesque is therefore doubly inexplicable. From where I am writing, I can, by turning a little, see the complete set of Bossuet’s books stacked up on the shelves familiar to me, and I have on this table, the Histoire des Variations, that I can open at my liking, assured of meeting in it the perfection of life from its first pages. Nevertheless, I would not know how to compare for its vivacity, the beautiful freshness of pleasure, that emotion already old to the emotion that the newly discovered chapters have given me.

A statue recently pulled out of the ground, even from an inferior age, touches and moves for an instant with more force than the most beautiful things known. Judge if chance brings to daylight a masterpiece from Phidias' school! We have here a masterpiece from Phidias himself. Space is lacking to justify my sentiment in regard to the profound and subtle ideas exposed in the second tract of the États d’oraison. But I am sure that we will tremble at the simple sound of the eloquence, its rhythm, and its performance.

In the chapter \textit{Que la foi se perd dans l'incompréhensibilité de Dieu} (``Faith is lost in the incomprehensibility of God''):

\begin{quotex}
It is a doctrine common among the holy Fathers and among theologians, mystical as well as scholastic, that we know God by negation rather than by affirmation; that is to say, that we know well what he is not, but we do not know what he is: so, in order to reach the lofty knowledge of the Divine Being, it is necessary to reject all our ideas and form our consciousness a little closer to the way we form a statue, by removing, from the marble from which it is formed, first one piece and then another, because it is the analogy which they use following the author of the Divine Names. So, they say, in order to know God, it is necessary to deny, in a certain sense, everything that we think of Him and everything that we say about Him, not as false, for it would be an impiety and an atheism to deny that God is holy, eternal, almighty, the ``I Am'', Father, Son, and Holy Spirit, and, in these three persons, one God: we don’t deny, therefore, these things as false, not pleasing to God! But we reject them in some fashion as still poorly proportionate and suitable to the immense perfection of the Divine Being: so that whatever effort that we make to know Him well, however sublime are the ideas that are presented to our spirits or the thoughts that we try to form from then, we deny that they are equal to his high and impenetrable majesty.
\end{quotex}  

In the chapter \textit{Theology and contemplation of saint Augustine}:

\begin{quotex}
But among the Fathers where this divine theology seems to me the most highly expressed, it is St. Augustine where he says to God: “I sought you outside, and you were within, I found you there without your entering by any of my senses. You didn’t come there with colours or exquisite tastes; you didn’t run there with odours or with pleasant sounds and chants. If you are a light, you are a light without a cloud, without sunset, and without body: you are nothing that I see, that I touch, that I sense, that I imagine in my mind, of that which I am: for my spirit, my mind, which is what I find most excellent in myself, learns, unlearns, forgets, loves certain things and then is disgusted by them; and God is not all that because He never changes. In rejecting, then, all those things and all those sense images, and opening only the eyes of the soul, it sees in itself, without any form, a justice which judges it and which it judges not, but by which it judges everything; a truth that slips away, for as little can one come near the sense of it as to touch it.
\end{quotex}

That is only a translation, a summary, a paraphrase; and genius breaks out in each word. Bossuet, moreover, excels in summaries. And he is perhaps the only classical Christian among whom the crudest citations of the Bible and of the Fathers, even though extremely numerous, do not subsequently result in a blemish in style. He establishes and transforms everything. His example can let us judge the value of modern theories of invention and originality.

But let us cite once more, a single sentence that ends on a charming inflexion: Faith makes it known to the soul that it is made in the image and likeness of God; this links two things: one, that it is pleased to gather its thoughts to contemplate in them the image of God; the other, that is can not stop in itself, but it is gently attracted and elevated to God, of whom it is the beautiful and living image. Yet, it is a mistake to blame oneself for admiring similar things in Bossuet. His spirit is too far above the sentences thus sketched out! But strength consists in finding the whole in its least objects and, thereby, according to a view loved by Hugo, even the small appears great.


\flright{From ``Barbarie et Poésie'', originally published in \textit{Revue Encyclopédique}, October 31, 1896.}